\documentclass[11pt,a4paper]{article}
% Shared preamble for RuPhO English problem statements (match T1 2026 style).
% Use: \input{latex/rupho_en_preamble.tex} after \documentclass.
\usepackage[margin=1in]{geometry}
\usepackage{amsmath,amssymb}
\usepackage{graphicx}
\usepackage{float}
\usepackage{enumitem}
\usepackage{microtype}
\usepackage{caption}
\usepackage{tabularx}
\usepackage{hyperref}

\captionsetup{font=small,labelfont=bf,skip=6pt}

\setlist[itemize]{leftmargin=1.2cm}
\setlist[enumerate]{leftmargin=1.2cm}

% One outer box per subproblem: [id | statement | points] (no inner rules)
\newcommand{\problembox}[3]{%
  \par\addvspace{0.42em}%
  \noindent
  \setlength{\fboxsep}{7.5pt}%
  \setlength{\fboxrule}{0.95pt}%
  \fbox{%
    \begin{minipage}{\dimexpr\linewidth-2\fboxsep-2\fboxrule\relax}
    \begin{tabularx}{\linewidth}{@{}>{\raggedright\arraybackslash\bfseries}p{2.1em} @{\hspace{0.9em}} >{\raggedright\arraybackslash}X @{\hspace{0.9em}} >{\raggedleft\arraybackslash\bfseries}p{2.4em}@{}}
    #1 & #2 & #3
    \end{tabularx}
    \end{minipage}%
  }%
  \par\addvspace{0.32em}%
}

\title{\textbf{Let's talk about the weather}}
\author{\normalsize X19 (2019) --- English translation}
\date{}
\begin{document}
\maketitle

The map below shows isobars (for the same altitude, approximately at sea level). Assume that the isobar diagram is stationary (changing slowly over time).

\begin{figure}[H]
    \centering
    \includegraphics[width=0.72\textwidth]{T5_EN_assets/fig1.jpg}
    \caption{Figure 1.}
\end{figure}

\problembox{A1}{Sketch the wind direction at points $A$ and $B$.}{1.00}

\problembox{A2}{Estimate the wind speed $v_A$ at point $A$. Consider the isobars to be straight lines around this point. Air density $\rho=1{.}23~\text{kg}/\text{m}^3$.}{2.00}

\problembox{A3}{Estimate the wind speed $v_B$ at point $B$.}{3.00}

Note: When a body of mass $m$ (for example, a layer of air) moves with speed $v$ in a rotating frame of reference (rotating with angular speed $\Omega$), it is acted upon by a certain fictitious force, which is called the Coriolis force: $$F_C=2m\Omega v\sin\varphi$$

\begin{figure}[H]
    \centering
    \includegraphics[width=0.72\textwidth]{T5_EN_assets/fig2.jpg}
    \caption{Figure 2.}
\end{figure}

\vspace{1em}
\noindent\textit{Source:} \url{https://pho.rs/p/3056}

\end{document}